\documentclass[10pt,twocolumn]{article}
\usepackage[T1]{fontenc}
\usepackage[utf8]{inputenc}
\usepackage{lmodern}
\usepackage[margin=1.8cm,top=2.4cm,bottom=2cm,columnsep=0.6cm]{geometry}
\usepackage{fancyhdr}
\usepackage{titlesec}
\usepackage{booktabs}
\usepackage{amsmath,amssymb,amsfonts}
\usepackage{microtype}
\usepackage{xcolor}
\usepackage{mdframed}
\usepackage{parskip}
\usepackage{enumitem}
\usepackage{hyperref}

\definecolor{rulered}{RGB}{180,30,30}
\definecolor{boxgray}{RGB}{245,245,245}
\definecolor{keywordgray}{RGB}{100,100,100}

\pagestyle{fancy}
\fancyhf{}
\fancyhead[L]{\textsc{\small Claude Mag}}
\fancyhead[C]{\small\textit{Machine Learning Research Digest}}
\fancyhead[R]{\small 2026-11-25}
\fancyfoot[C]{\small\thepage}
\renewcommand{\headrulewidth}{0.8pt}

\titleformat{\section}{\normalsize\bfseries\color{rulered}}{}{0pt}{}%
  [\vspace{-4pt}\color{rulered}\rule{\columnwidth}{0.4pt}]
\titlespacing{\section}{0pt}{8pt}{4pt}

\hypersetup{colorlinks=true,urlcolor=rulered,linkcolor=black}

\begin{document}
\setlength{\parindent}{0pt}
\setlength{\parskip}{4pt}

{\small\color{keywordgray}\textbf{Keywords:}
  SFT-then-RL \textbullet{} Mixed-Policy Training \textbullet{} LLM Reasoning \textbullet{} Optimizer Bug \textbullet{} Reproducibility}\par\vspace{4pt}

{\LARGE\bfseries\raggedright SFT-then-RL Outperforms Mixed-Policy Methods for LLM Reasoning}\par\vspace{4pt}
{\large\itshape\raggedright Two bugs in the standard baseline were suppressing performance; after fixes, the sequential pipeline beats every published mixed-policy method by up to 22 points}\par\vspace{6pt}
{\small\textbf{By} Alexis Limozin, Eduard Durech, Torsten Hoefler, Imanol Schlag, Valentina Pyatkin}\par\vspace{2pt}
{\small\color{keywordgray}arXiv:2604.23747 \textbullet{} cs.LG, cs.CL \textbullet{} NeurIPS 2026}
\par\vspace{2pt}\noindent{\color{rulered}\rule{\columnwidth}{0.6pt}}\vspace{6pt}

Limozin and colleagues identify an optimizer bug and a loss aggregation bug in the standard SFT-then-RL baseline used in every published mixed-policy comparison.
After fixing both, sequential SFT-then-RL outperforms LUFFY, ReLIFT, SRFT, and two other mixed-policy methods by up to $+22.2$ points on math benchmarks---calling into question much of the case for interleaved training.

\begin{mdframed}[backgroundcolor=boxgray,linecolor=rulered,linewidth=1pt,
    innerleftmargin=6pt,innerrightmargin=6pt,
    innertopmargin=6pt,innerbottommargin=6pt]
\textbf{\small AT A GLANCE}\par\vspace{4pt}
\begin{description}[leftmargin=0pt,labelsep=4pt,itemsep=2pt,parsep=0pt,
    font=\small\bfseries]
  \item[Problem:] Published comparisons of mixed-policy methods against sequential SFT-then-RL used a buggy baseline, inflating the apparent advantage of mixed approaches.
  \item[Why it matters:] A large body of recent work has built increasingly complex interleaved training algorithms to close a gap that may not exist.
  \item[Method:] Identify and correct two bugs (optimizer state reset, loss aggregation); re-evaluate five published mixed-policy methods against the fixed baseline.
  \item[Result:] +3.8 pts (Qwen2.5-Math-7B); +22.2 pts (Llama-3.1-8B) over best mixed-policy method; 50-step truncated variant beats mixed approaches at 0.3$\times$ FLOPs.
  \item[Evidence:] OpenR1-Math benchmark; Qwen2.5-Math-7B and Llama-3.1-8B; code at github.com/alek6kun/sft\_then\_rl.
\end{description}
\end{mdframed}

\section{The Big Picture}

RLVR training for LLM reasoning follows a standard recipe: supervised fine-tuning on correct reasoning chains, followed by policy gradient updates guided by a verifiable reward.
The simplest instantiation---sequential SFT-then-RL---was reported to perform worse than more complex ``mixed-policy'' methods that interleave SFT and RL data or updates within a single stage.
This motivated a wave of work designing ever-more-sophisticated interleaving algorithms.

This paper traces the apparent advantage of mixed-policy methods to implementation errors in the sequential baseline.
Once those errors are corrected, the picture reverses: sequential SFT-then-RL is both simpler and stronger.

\section{Why Existing Approaches Fall Short}

Mixed-policy methods were designed to solve a problem they did not need to solve.
The leading hypothesis in prior work was that the two-stage sequential pipeline fails because the RL stage is initialized from a checkpoint that has ``overfit'' to SFT data, limiting its ability to explore new strategies.
Mixed-policy methods addressed this by keeping the policy on-distribution throughout training.

The alternative hypothesis---that the sequential baseline was simply misconfigured---was not tested.
This paper tests it, finding that the misconfiguration explains the entire gap.

\section{The Two Bugs}

\subsection*{Bug 1: Stale Optimizer State}

When the pipeline transitions from SFT to RL, the optimizer should be reset.
SFT uses cross-entropy loss; RL uses a policy gradient loss with advantages as coefficients.
The gradient statistics accumulated by AdamW during SFT reflect a fundamentally different loss landscape.
Using those statistics to initialize the RL phase produces miscalibrated step sizes and distorted momentum, destabilizing the first few hundred RL updates.

The fix: zero the optimizer's first and second moment estimates at the SFT-to-RL transition and restart the learning rate schedule.

\subsection*{Bug 2: Loss Aggregation}

Policy gradient loss must normalize over sequence length to weight all training examples equally, regardless of how long their responses are.
The buggy implementation used batch-level mean reduction (dividing total loss by total tokens in a minibatch), rather than sequence-level mean reduction (averaging per-sequence means across the minibatch):

\begin{equation}
\text{Buggy:}\quad \mathcal{L} = \frac{\sum_{i,t} A_{i,t} \log \pi_\theta(y_{i,t})}{\sum_i |y_i|}
\end{equation}

\begin{equation}
\text{Fixed:}\quad \mathcal{L} = \frac{1}{|B|}\sum_i \frac{1}{|y_i|}\sum_t A_{i,t} \log \pi_\theta(y_{i,t})
\end{equation}

Batch-level reduction down-weights long responses: a 400-token reasoning chain contributes the same total loss as four 100-token chains, but four times less per sequence.
Since RLVR is intended to encourage long, correct reasoning chains, this bug systematically penalizes the behavior being trained.

\section{Learning or Inference Procedure}

The fixed SFT-then-RL pipeline:
\begin{enumerate}[leftmargin=1.5em,itemsep=1pt]
  \item Train with cross-entropy on SFT dataset to convergence.
  \item \textbf{Reset optimizer state.} Restart learning rate schedule.
  \item Train with GRPO using correctness reward and \textbf{sequence-level loss normalization}.
  \item Evaluate on held-out math benchmarks.
\end{enumerate}

The truncated variant runs only 50 RL steps in stage~3.

\section{Experimental Findings}

\begin{center}
\small
\begin{tabular}{lrr}
\toprule
\textbf{Method} & \textbf{Qwen-7B} & \textbf{Llama-8B} \\
\midrule
LUFFY & 61.4 & 46.3 \\
ReLIFT & 62.1 & 47.1 \\
SRFT & 62.8 & 48.3 \\
SFT-then-RL (buggy) & 59.2 & 46.1 \\
SFT-then-RL (fixed) & \textbf{66.6} & \textbf{70.5} \\
SFT-then-RL-50 (fixed) & 65.1 & 68.4 \\
\bottomrule
\end{tabular}
\end{center}

The 22.2-point gap on Llama-3.1-8B reflects that model's higher sensitivity to the optimizer bug---its SFT checkpoint is farther from the RL-optimal initialization than Qwen2.5-Math-7B, so stale optimizer state causes larger initial instability.

The 50-step truncated variant achieves 97.7\% of the full fixed pipeline's performance at 0.3$\times$ FLOPs.

\section{Ablations and Interpretation}

\begin{itemize}[leftmargin=1em,itemsep=2pt]
  \item \textbf{Optimizer bug alone:} Fixing only the optimizer reset improves performance by +4.1 points on Qwen-7B.
  \item \textbf{Loss bug alone:} Fixing only the normalization improves performance by +2.8 points.
  \item \textbf{Both bugs:} Combined fix produces +7.4 points---confirming super-additive interaction.
  \item \textbf{Generalization:} The fixed pipeline generalizes beyond the OpenR1-Math training distribution: on AIME 2026, the fixed pipeline outperforms the best mixed-policy method by +6.2 points.
  \item \textbf{Stability:} Training loss variance in the first 100 RL steps is $3.4\times$ lower with the optimizer reset, confirming the stabilization mechanism.
\end{itemize}

\section{What the Guarantee Says}

The paper does not claim that mixed-policy methods have no merit in principle.
It claims that the specific comparisons made in prior work were confounded by implementation errors.
Mixed-policy methods may still offer advantages in settings with very limited SFT data, non-stationary reward functions, or online environments where the SFT checkpoint ages quickly.
But the burden of proof is now on mixed-policy methods to demonstrate advantages over a correctly implemented sequential baseline.

\section*{Reference}
A.\ Limozin, E.\ Durech, T.\ Hoefler, I.\ Schlag, V.\ Pyatkin.
\textbf{SFT-then-RL Outperforms Mixed-Policy Methods for LLM Reasoning.}
arXiv:2604.23747. cs.LG, cs.CL. NeurIPS 2026.
\url{https://arxiv.org/abs/2604.23747}

\end{document}
